% =====================================================================
% Paper B
% Direction-averaged whole-body absorbed-power identity:
% closed-form dosimetry from 100 MHz to 100 GHz,
% validated against three independent measurement campaigns.
% Author: Robin Wydaeghe
% Target: IEEE Transactions on Antennas and Propagation
% =====================================================================
\documentclass[journal,twocolumn,10pt]{IEEEtran}

\usepackage[utf8]{inputenc}
\usepackage[T1]{fontenc}
\usepackage{lmodern}
\usepackage{microtype}

\usepackage{amsmath,amssymb,amsthm,mathtools}
\usepackage{bm}

\usepackage{graphicx}
\graphicspath{{../../theory/figures/}{../../validation/}{figures/}}
\usepackage{booktabs}
\usepackage{array}
\usepackage{caption}
\usepackage{subcaption}
\usepackage{cite}

\usepackage{xcolor}
\usepackage[colorlinks=true,linkcolor=blue,citecolor=blue,urlcolor=blue]{hyperref}
\usepackage[capitalize]{cleveref}

% Theorem environments
\theoremstyle{plain}
\newtheorem{theorem}{Theorem}
\newtheorem{proposition}[theorem]{Proposition}
\newtheorem{corollary}[theorem]{Corollary}
\theoremstyle{remark}
\newtheorem{remark}[theorem]{Remark}

\crefname{proposition}{Proposition}{Propositions}
\Crefname{proposition}{Proposition}{Propositions}
\crefname{corollary}{Corollary}{Corollaries}
\Crefname{corollary}{Corollary}{Corollaries}
\crefname{remark}{Remark}{Remarks}
\Crefname{remark}{Remark}{Remarks}

% Math macros
\newcommand{\khat}{\hat{\bm{k}}}
\newcommand{\nhat}{\hat{\bm{n}}}
\newcommand{\rr}{\mathbf{r}}
\newcommand{\Sinc}{S_{\mathrm{inc}}}
\newcommand{\Sab}{S_{\mathrm{ab}}}
\newcommand{\Aab}{A_{\mathrm{ab}}}
\newcommand{\Aperp}{A_\perp}
\newcommand{\Tavg}{T_{\mathrm{avg}}}
\newcommand{\Tbar}{\bar{T}}
\newcommand{\ntilde}{\tilde{n}}
\newcommand{\diff}{\mathrm{d}}
\newcommand{\pospart}[1]{\left(#1\right)_{+}}
\newcommand{\Vis}{V}
\DeclareMathOperator{\ReLU}{ReLU}
\DeclareMathOperator{\RE}{Re}

\begin{document}

\title{Direction-averaged whole-body absorbed power: a closed-form identity from 100\,MHz to 100\,GHz}

\author{Robin~Wydaeghe%
  \thanks{R. Wydaeghe is with Ghent University and imec, Ghent, Belgium
  (e-mail: robin.wydaeghe@ugent.be).}}

\markboth{IEEE Transactions on Antennas and Propagation, Vol.~XX, No.~X, Month~2026}%
{Wydaeghe: Direction-averaged whole-body absorbed power}

\maketitle

\begin{abstract}
The direction-averaged whole-body absorbed power on a human body
exposed to a richly multipath or reverberant electromagnetic field is
$\langle P_{\mathrm{abs}} \rangle = \Sinc\,\Tbar(f)\,\Aab/4$, exact at
any frequency where the body is opaque, that is, above
$\sim 100\,$MHz. Here $\Tbar(f)$ is the flux-weighted Fresnel
transmission of skin, computed from the IT'IS dielectric tables by one
quadrature, and $\Aab = \int \eta\,\diff A$ is the body surface area
weighted pointwise by an ambient-occlusion factor $\eta$. The factor
$\eta$ coincides with the self-shadowing ratio $\gamma_s$ used by
Flintoft and the non-convexity correction implied by Zhang. The
classical Cauchy projected-area formula of 1841 is the convex limit
$\eta \equiv 1$, $\Tbar \equiv T_0$. The dip near 3\,GHz observed in
volunteer reverberation-chamber data is a Fabry-P\'erot resonance in
the subcutaneous fat layer, modelled by a three-layer transfer-matrix
construction. The identity matches Sim4Life FDTD on the Thelonious anatomical
phantom to 1.2\,\% at 5.8\,GHz with no fitted parameters, and
reproduces every plateau, dip, and BSA scaling reported in 168
volunteers across the Bamba 2014, Flintoft 2014, and Zhang 2017
campaigns, plus the numerical campaigns of Diao 2024 (28\,GHz) and
Kodera 2024 (1--100\,GHz). Whole-body SAR compliance reduces to three
precomputed scalars: $\Tbar$, $\Aab$, and the body mass $m$.
Reverberation-chamber dosimetry reduces to one closed-form evaluation.
\end{abstract}

\begin{IEEEkeywords}
Absorption cross-section, ambient occlusion, Cauchy formula,
electromagnetic dosimetry, Fabry-P\'erot resonance, Fresnel
transmission, ICNIRP compliance, reverberation chamber, specific
absorption rate, whole-body absorbed power.
\end{IEEEkeywords}

\IEEEpeerreviewmaketitle


\section{Introduction}\label{sec:intro}
\IEEEPARstart{F}{or} a quarter of a century the bioelectromagnetics
community has measured the same physical quantity in five different
notations. In a reverberation chamber, the whole-body absorbed power
of a human volunteer divided by the incident power density is the
absorption cross-section $\sigma_a$ in m$^2$. Different groups have
normalised $\sigma_a$ in different ways. Bamba
et al.~\cite{Bamba2014} divide by phantom mass to get a
direction-averaged whole-body SAR and fit an empirical efficiency
$\eta(f) \approx 0.50$ to $0.56$ across 1.45 to 5.8\,GHz. Flintoft
et al.~\cite{Flintoft2014} divide by $\gamma_s\,\mathrm{BSA}/4$ to
get a normalised cross-section $\langle Q^a\rangle$, where
$\gamma_s$ is the body self-shadowing factor and BSA is the body
surface area. They report $\langle Q^a\rangle \approx 0.41$ from 7 to
11\,GHz on 60 volunteers. Zhang~\cite{Zhang2017thesis,
ZhangRobinson2020} divides $\sigma_a$ by $\mathrm{BSA}/4$ to get a
coefficient $\xi$ of 0.45 to 0.65 above 6\,GHz on 48 subjects.
Independently, Kodera et al.~\cite{Kodera2024} fit a
transmission coefficient $T_{\mathrm{tr}}$ such that
$\mathrm{SAR}_{\mathrm{wb}} = T_{\mathrm{tr}}\,A_\perp\,\Sinc/W$
matches FDTD to 5\,\% from 10 to 100\,GHz on parametric and anatomical
phantoms. Diao et al.~\cite{Diao2024} report $T = 0.52$ at
28\,GHz from full anatomical FDTD. None of these works derives the
limit from Maxwell's equations, none extends the result to non-convex
bodies in closed form, and none explains the persistent dip in
$\langle Q^a\rangle(f)$ and $\xi(f)$ near 3\,GHz. Flintoft attributes
the dip qualitatively to ``reflections between layers of
tissue''~\cite[p.~3306]{Flintoft2014}. Zhang derives a planar
skin/fat/muscle transfer-matrix model that predicts the
dip~\cite[Sec.~2.2]{Zhang2017thesis} but does not embed it in a
body-surface average. Kodera et al.~\cite{Kodera2024} report the
most extensive numerical comparison to date, fitting a study-specific
$T_{\mathrm{tr}}$ to each phantom and frequency. They ask, but do not
answer, what closed-form $T_{\mathrm{tr}}(f)$ all phantoms converge
to at high frequency.

The direction-averaged whole-body absorbed power on any body opaque
at the wavelength is
\begin{equation}\label{eq:cauchy-id}
  \langle P_{\mathrm{abs}} \rangle = \Sinc\,\Tbar(f)\,\Aab/4,
\end{equation}
exact above $\sim 100\,$MHz. Here $\Tbar(f)$ is the flux-weighted
Fresnel transmission of skin, obtained by one quadrature against the
IT'IS v5.0 Cole-Cole dielectric model, and $\Aab = \int_\Sigma
\eta(\rr)\,\diff A$ is the surface area weighted pointwise by the
exposure fraction $\eta$. The classical Cauchy projected-area formula
of 1841~\cite{Cauchy1841} is the special case of a convex body with
$\eta \equiv 1$ and $\Tbar \equiv T_0$, the normal-incidence Fresnel
transmission. The body self-shadowing factor $\gamma_s$ that Flintoft
defines verbatim as ``the proportion of total surface area of the
body that is illuminated by the reverberant
field''~\cite[p.~3301]{Flintoft2014} is the area-weighted mean of
$\eta$. The exposure fraction $\eta(\rr)$ is mathematically identical
to the ambient occlusion factor of computer
graphics~\cite{Zhukov1998,Landis2002,AkenineMoller2018}, a real-time
GPU primitive. For the Thelonious anatomical phantom, the
ambient-occlusion solver in AEGIS~\cite{AEGIS2026} returns
$\Aab/A = 0.865$, agreeing with Flintoft's geometric estimate from
Tomita's surface-area data~\cite{Tomita1999} to within the upper end
of his $0.75$--$0.85$ band. The dip near 3\,GHz is the Fabry-P\'erot
resonance in the subcutaneous fat layer, predicted by a three-layer
transfer matrix and consistent with Flintoft's correlation slope of
$\langle Q^a\rangle$ against subcutaneous fat thickness, which peaks
at 3\,GHz and falls by a factor of two at 7\,GHz.
Equation~\eqref{eq:cauchy-id} matches Sim4Life FDTD on the Thelonious
phantom to 1.2\,\% at 5.8\,GHz, matches every plateau and dip in 168
volunteers across the three reverberation-chamber campaigns, and
matches the numerical campaigns of Diao 2024 and Kodera 2024 to
within 5\,\%. Whole-body SAR compliance under the ICNIRP 2020
guidelines~\cite{ICNIRP2020} reduces to three precomputed scalars
($\Tbar$, $\Aab$, $m$). Reverberation-chamber dosimetry on small
animals or human-body phantoms reduces to one closed-form evaluation.
The paper proceeds as follows. Section~\ref{sec:cauchy} states and
proves the generalised Cauchy formula for non-convex bodies.
Section~\ref{sec:Tbar} derives the flux-weighted transmission $\Tbar$
from the Fresnel coefficients and tabulates it from 0.3 to 100\,GHz.
Section~\ref{sec:fp} treats the Fabry-P\'erot dip via the layered
tissue model. Section~\ref{sec:compliance} reduces ICNIRP whole-body
SAR compliance to closed form and gives the anthropometric scaling.
Section~\ref{sec:waterfall} compares the identity with the
reverberation-chamber and FDTD literature.
Section~\ref{sec:fdtd} reports the direct FDTD validation.
Sections~\ref{sec:reverb} and~\ref{sec:concl} cover reverberation-
chamber implications and conclude.


\section{Generalised Cauchy formula}\label{sec:cauchy}

The classical projected-area formula of Cauchy~\cite{Cauchy1841}
states that for a convex body with surface area $A$, the
direction-averaged silhouette area equals $A/4$. The generalisation
below extends the formula to non-convex bodies and replaces the
silhouette area with an absorption area that accounts for self-
shadowing.

\subsection{Local absorption law}

For a plane wave with power density $\Sinc$ and direction $\khat$
incident on biological tissue, the absorbed power surface density
satisfies
\begin{equation}\label{eq:Sab-local}
  \Sab(\rr) \approx \Sinc\,T_0(f)\,\Vis(\rr,\khat)\,
  \pospart{\nhat(\rr)\cdot(-\khat)},
\end{equation}
where $T_0(f)$ is the normal-incidence Fresnel transmission of skin,
$\nhat(\rr)$ is the outward unit normal at $\rr$, $(\cdot)_+ \equiv
\max(0,\cdot)$ is the positive-part operator that gates back-facing
elements to zero, and $\Vis(\rr,\khat) \in \{0,1\}$ is the binary
visibility, equal to one when $\rr$ is reached by an unobstructed ray
from direction $\khat$ and zero when $\rr$ is shadowed by another
body part. In machine-learning notation $(\cdot)_+$ is the ReLU
activation. The constant-$T_0$ approximation is exact at the
pseudo-Brewster crossover near 40\,GHz and matches the full
angle-resolved Fresnel result to 5\,\% root-mean-square across 0.3 to
100\,GHz on skin. The derivation uses opacity at the wavelength to
truncate the body integral to the lit surface, with the Mie-analysis
bridge for body-scale spheres in Bohren and Huffman~\cite{BohrenHuffman1983}.
For a convex body $\Vis \equiv 1$ and the formula reduces to the
classical $\Sinc T_0 \pospart{\nhat\cdot(-\khat)}$. Integrating
\eqref{eq:Sab-local} over the illuminated surface of a convex body
gives
\begin{equation}\label{eq:Pabs-convex}
  P_{\mathrm{abs}}(\khat) = \Sinc\,T_0\int_{\Sigma_+}\cos\theta\,\diff A
  = \Sinc\,T_0\,\Aperp(\khat),
\end{equation}
with $\Aperp(\khat) = \int_{\Sigma_+}\cos\theta\,\diff A$ the
projected (shadow) area perpendicular to $\khat$. Kodera
et al.~\cite{Kodera2024} write the same identity as
$\mathrm{SAR}_{\mathrm{wb}} = T_{\mathrm{tr}}\,A_\perp\,\Sinc/W$, with
$T_{\mathrm{tr}}$ extracted from a one-dimensional FDTD slab model.
We identify $T_{\mathrm{tr}}$ with $T_0$ in the convex limit.

\subsection{Self-shadowing and ambient occlusion}\label{sec:self-shadow}

The human body is not convex. Concavities such as the armpits, the
gap between the legs, and the neck region cause one part of the body
to shadow another. The binary visibility $\Vis(\rr,\khat) \in
\{0,1\}$ in \eqref{eq:Sab-local} is the ambient-occlusion primitive
of computer graphics introduced by Zhukov et al.~\cite{Zhukov1998}
and brought into production rendering by Landis~\cite{Landis2002}.
Modern GPUs evaluate $\Vis(\rr,\khat)$ at interactive frame
rates~\cite{AkenineMoller2018}.

\begin{remark}[Exposure fraction]\label{def:eta}
The exposure fraction at a surface point $\rr$ is the cosine-weighted
fraction of the upper hemisphere from which $\rr$ is unobstructed,
\begin{equation}\label{eq:eta-def}
  \eta(\rr) = \frac{1}{\pi}\int_{S^2}
  \pospart{\nhat(\rr)\cdot(-\khat)}\,\Vis(\rr,\khat)\,\diff\Omega(\khat).
\end{equation}
For convex bodies $\Vis \equiv 1$ and $\eta \equiv 1$. For non-convex
bodies $\eta \in [0,1]$.
\end{remark}

The identification of $\eta$ with Flintoft's self-shadowing factor
$\gamma_s$ is anticipated by Flintoft et al., who define the
self-shadowing factor as ``the proportion of total surface area of
the body that is illuminated by the reverberant
field''~\cite[p.~3301]{Flintoft2014} and estimate $\gamma_s = 0.75$
to $0.85$ from Tomita's surface-area data~\cite{Tomita1999}. The Flintoft estimate is geometric and
posture-dependent. The present construction is computational and
posture-resolved. For the Thelonious phantom, AEGIS's ambient-
occlusion solver returns the area-weighted mean
$\bar{\eta} = \Aab/A = 0.865$, sitting near the upper end of the
Flintoft band. \Cref{fig:eta-phantom} renders $\eta(\rr)$ on the
phantom surface.

\begin{figure}[!t]
  \centering
  \includegraphics[width=\columnwidth]{eta_phantom_paperB.pdf}
  \caption{Exposure fraction $\eta(\rr) \in [0,1]$ on the Thelonious
  anatomical phantom from the cosine-weighted hemispheric integral
  \eqref{eq:eta-def}, front and side views. Most of the body sits
  near $\eta \approx 1$ (limbs and torso facing outward), with
  reductions concentrated in concavities: the medial sides of the
  legs and arms, the armpits, and the underside of the chin and
  feet. The area-weighted mean is reported below the panels and
  agrees with Flintoft's geometric estimate from Tomita's
  surface-area data ($0.75$ to $0.85$~\cite{Tomita1999}). The map
  is computed on the same mesh used in the FDTD validation in
  \cref{sec:fdtd}.}
  \label{fig:eta-phantom}
\end{figure}

\subsection{Statement and proof}

\begin{theorem}[Generalised Cauchy formula]\label{thm:cauchy}
Let a body $\Sigma$ have surface area $A$, exposure fraction
$\eta(\rr)$, and absorption area
$\Aab \equiv \int_\Sigma \eta(\rr)\,\diff A$. Under isotropic,
unpolarised plane-wave illumination of intensity $\Sinc$ on tissue
with normal-incidence transmission $T_0$, the direction-averaged
whole-body absorbed power is
\begin{equation}\label{eq:cauchy}
  \langle P_{\mathrm{abs}} \rangle = \Sinc\,T_0\,\Aab/4.
\end{equation}
\end{theorem}

\begin{proof}
Apply Fubini's theorem to exchange the surface and direction
integrals. The local law \eqref{eq:Sab-local} gives
$\Sab(\rr,\khat) = \Sinc\,T_0\,\Vis(\rr,\khat)\,
\pospart{\nhat\cdot(-\khat)}$. The direction average of the
integrand is
\[
  \frac{1}{4\pi}\int_{S^2}
  \Sinc\,T_0\,\Vis(\rr,\khat)\,\pospart{\nhat\cdot(-\khat)}\,\diff\Omega
  = \frac{\Sinc\,T_0}{4}\,\eta(\rr),
\]
using the definition of $\eta$ and the identity
$\int_{S^2}\pospart{\nhat\cdot(-\khat)}\,\diff\Omega = \pi$ for any unit
$\nhat$. Integration over $\Sigma$ gives~\eqref{eq:cauchy}.
\end{proof}

The classical Cauchy formula $\langle\Aperp\rangle = A/4$ is the
special case $\eta \equiv 1$, valid for any convex body. The
absorption area $\Aab$ collects all geometric complexity of self-
shadowing into a single scalar.

\begin{remark}[Convex-hull energy bound]\label{rem:hull}
Let $A_{\mathrm{CH}}$ be the surface area of the convex hull of the
body. Energy conservation under isotropic illumination implies
$\langle P_{\mathrm{abs}} \rangle \leq \Sinc\,A_{\mathrm{CH}}/4$,
because the power entering the convex hull bounds the absorbed power.
For the Thelonious phantom $A_{\mathrm{CH}}/A \approx 1.20$, so the
hull bound brackets the true absorbed power within a few percent.
\end{remark}


\section{Skin transmission across frequency}\label{sec:Tbar}

The constant-$T_0$ approximation in~\eqref{eq:Sab-local} is accurate
to 5\,\% root-mean-square across 0.3 to 100\,GHz, but it is not
exact. Replacing $T_0$ in~\eqref{eq:cauchy} with a flux-weighted
transmission $\Tbar(f)$ removes the approximation and gives a
formula valid at any frequency where the body is opaque.

\begin{remark}[Flux-weighted transmission]\label{def:Tbar}
The flux-weighted Fresnel transmission is
\begin{equation}\label{eq:Tbar-def}
  \Tbar(f) \equiv 2\int_0^1 \Tavg(\mu,f)\,\mu\,\diff\mu,
\end{equation}
where $\Tavg(\mu,f) = \tfrac{1}{2}(T_s + T_p)$ is the unpolarised
Fresnel transmission and $\mu = \cos\theta$ is the incidence cosine.
$\Tbar$ is the fraction of incident power absorbed by a convex body
under isotropic, unpolarised illumination.
\end{remark}

\begin{theorem}[Exact direction-averaged absorbed power]\label{thm:exact}
For any body opaque at the wavelength,
\begin{equation}\label{eq:cauchy-exact}
  \langle P_{\mathrm{abs}} \rangle = \Sinc\,\Tbar(f)\,\Aab/4.
\end{equation}
\end{theorem}

\begin{proof}
Replace $T_0$ in the proof of \cref{thm:cauchy} with the angle-
dependent $\Tavg(\theta)$. The cosine-weighted angular integral
collapses to $\Tbar$ via~\eqref{eq:Tbar-def}.
\end{proof}

Equation~\eqref{eq:cauchy-exact} requires only electromagnetic
opacity, a condition met above $\sim 100\,$MHz on a human body, with
no pseudo-Brewster approximation invoked. \cref{tab:Tbar} lists $T_0$,
$\Tbar$, and the ratio $R = T_0/\Tbar$ for skin from 0.3 to 100\,GHz,
computed from the IT'IS v5.0 Cole-Cole dielectric
model~\cite{ITISv5,Gabriel1996}. The crossover $R = 1$ falls at
40.4\,GHz, where the constant-$T_0$ form is exact to four significant
figures. Below this frequency the simplified formula with $T_0$
underestimates absorbed power and is therefore conservative for
compliance. Above, it overestimates by at most 3.5\,\% at 100\,GHz.

\begin{table}[!t]
\centering
\caption{Normal-incidence transmission $T_0$, flux-weighted
transmission $\Tbar$, and their ratio $R = T_0/\Tbar$ for skin
(IT'IS v5.0 Cole-Cole model).}\label{tab:Tbar}
\begin{tabular}{rcccc}
\toprule
$f$\,(GHz) & $|\ntilde|$ & $T_0$ & $\Tbar$ & $R$ \\
\midrule
0.3 & 7.93 & 0.381 & 0.406 & 0.938 \\
0.9 & 6.70 & 0.445 & 0.465 & 0.957 \\
2.4 & 6.29 & 0.470 & 0.488 & 0.964 \\
6.0 & 6.07 & 0.481 & 0.497 & 0.968 \\
10  & 5.87 & 0.489 & 0.504 & 0.971 \\
28  & 4.84 & 0.536 & 0.543 & 0.988 \\
40  & 4.30 & 0.570 & 0.571 & 1.000 \\
60  & 3.68 & 0.622 & 0.613 & 1.015 \\
100 & 3.01 & 0.701 & 0.677 & 1.035 \\
\bottomrule
\end{tabular}
\end{table}

The reverberation-chamber community has been measuring $\Tbar$
directly. Bamba's empirical efficiency $\eta(f)$ for diffuse-field
exposure tracks $\Tbar(f)$ within 5 to 10\,\% across 1.45 to
5.8\,GHz~\cite{Bamba2014}. The systematic offset is consistent with
finite-curvature creeping-wave contributions on their ellipsoidal
phantoms (Section~\ref{sec:waterfall}). Flintoft's plateau
$\langle Q^a\rangle/\gamma_s = 0.47$ to $0.49$ at 7 to 11\,GHz matches
$\Tbar$ at the same frequencies to 2\,\%~\cite{Flintoft2014}. Zhang's
plateau $\xi = 0.45$ to $0.65$ above 6\,GHz brackets
$\Tbar\cdot\Aab/A$~\cite{Zhang2017thesis}.


\section{Fat-layer resonance below 5\,GHz}\label{sec:fp}

\Cref{thm:exact} reproduces the plateau values reported above 6\,GHz
across all three measurement campaigns. Below this frequency, both
Flintoft and Zhang observe a structured dip near 3\,GHz that the
homogeneous half-space model in Section~\ref{sec:Tbar} does not
capture. The dip is anatomical rather than instrumental: Flintoft's
negative correlation of $\langle Q^a\rangle$ with mean subcutaneous
fat thickness $d_{\mathrm{SF}}$ is steepest at 3\,GHz
($-0.0061\,\mathrm{mm}^{-1}$, $R^2 = 0.40$~\cite[Table~6]{Flintoft2014}),
with the slope falling to $-0.0030\,\mathrm{mm}^{-1}$ at 7 to 11\,GHz.

The mechanism is a Fabry-P\'erot resonance in the subcutaneous fat
layer. Below 6\,GHz the SAR penetration depth in fat exceeds 70\,mm,
against fat thicknesses of 2 to 20\,mm in the Flintoft
cohort~\cite[Table~1]{Flintoft2014}. The wave traverses the fat layer
with little attenuation and reflects from the fat-muscle interface.
Constructive interference enhances absorption, and destructive
interference suppresses it. A three-layer transfer-matrix model with
skin, fat, and a semi-infinite muscle half-space gives the layered
transmission
\begin{equation}\label{eq:T-lay}
  T_{\mathrm{lay}}(f, d_{\mathrm{SF}})
  = 1 - \bigl|\widetilde{\Gamma}_1(f, d_{\mathrm{SF}})\bigr|^2,
\end{equation}
where $\widetilde{\Gamma}_1$ is the generalised Fresnel reflection
coefficient at the air-skin interface, computed recursively from the
fat-muscle interface upward~\cite{Chew1995,BornWolf1999}. Replacing
$\Tbar$ with $T_{\mathrm{lay}}$ in~\eqref{eq:cauchy-exact} gives a
frequency-dependent direction-averaged absorbed power that includes
the fat-layer resonance. For a fat thickness of 10\,mm the model
predicts a 40\,\% enhancement above the homogeneous prediction at
0.9\,GHz (quarter-wave matching) and a 27\,\% reduction at 3.5\,GHz
(destructive interference).

Zhang derives the planar limit of this model in his
thesis~\cite[Sec.~2.2]{Zhang2017thesis} and observes the resonance
shift with fat thickness in his Figs.~2.7--2.8. He writes that
``the fat layer may act as a matching layer between skin and
muscle''~\cite[p.~21]{Zhang2017thesis}. The contribution here is to
embed the planar model in the body-surface integral via
$T_{\mathrm{lay}}$ and \cref{thm:exact}, which makes the resonance
compatible with a Cauchy-style direction average. Body-surface
averaging suppresses the oscillation by a factor of approximately
$A_{\mathrm{exposed,fat}}/A$ because different body regions carry
different fat thicknesses, with limbs near 2\,mm and abdomen near
30\,mm, and consequently different resonance frequencies. The
integrated dip is shallower than the single-thickness prediction but
sits at the same frequency. The local surface map
$\Sab(\rr) \approx \Sinc T_0 \Vis \pospart{\mu}$ loses pointwise meaning
below 6\,GHz, where the SAR penetration depth
exceeds the surface layer thickness. Total power remains valid via
$T_{\mathrm{lay}}$ throughout.


\section{Closed-form ICNIRP compliance}\label{sec:compliance}

The ICNIRP 2020 guidelines~\cite{ICNIRP2020} specify a whole-body
average SAR limit of 0.08\,W/kg for the general public. \Cref{thm:exact}
reduces compliance assessment to three precomputed scalars and
gives the anthropometric scaling in closed form.

\subsection{Compliance threshold}

Whole-body SAR is $\mathrm{SAR}_{\mathrm{wb}} = P_{\mathrm{abs}}/m$
for body mass $m$. The hemisphere bound on the absorption directivity
$D(\khat) = 4\Aperp(\khat)/\Aab \leq 2\,A_{\mathrm{CH}}/\Aab$ gives
the worst-case direction-resolved threshold
\begin{equation}\label{eq:Sinc-max}
  \Sinc < \frac{0.32\,m}{\Tbar\,\Aab\,D(\khat)},
\end{equation}
or, with $D \leq 2$,
\begin{equation}\label{eq:Sinc-max-worst}
  \Sinc < \frac{0.16\,m}{\Tbar\,\Aab}.
\end{equation}
Three quantities determine the threshold: the body mass $m$, the
absorption area $\Aab$ (equivalently the BSA times the self-shadowing
factor), and the tissue transmission $\Tbar$. None requires an FDTD
solve on the specific exposure scenario.

\subsection{Anthropometric scaling}

Body surface area follows the Du Bois formula~\cite{DuBois1916}
$A \approx 0.007184\,m^{0.425}\,h^{0.725}$ with mass in kg and height
in cm. The compliance threshold scales as $m/A \propto
\mathrm{BMI}^{0.575}\,h^{0.425}$. Heavier and taller individuals
tolerate higher incident power densities. \Cref{tab:anthro} evaluates
\eqref{eq:Sinc-max-worst} on a representative population at
28\,GHz with $T_0 = 0.536$ and $\gamma_s = 0.85$.

\begin{table}[!t]
\centering
\caption{Compliance threshold across the human population at 28\,GHz
with $T_0 = 0.536$ and $\gamma_s = 0.85$. The threshold varies by a
factor of about two between an infant and a large adult.}
\label{tab:anthro}
\small
\setlength{\tabcolsep}{3pt}
\begin{tabular}{lcccc}
\toprule
Body type & $m$\,(kg) & $h$\,(cm) & $m/A$ & $\Sinc^{\max}$ \\
& & & (kg/m$^2$) & (W/m$^2$) \\
\midrule
Infant (1\,yr) & 8   & 70  & 21.1 & 6.2 \\
Child (6\,yr)  & 20  & 110 & 25.8 & 7.6 \\
Adolescent     & 50  & 160 & 33.5 & 9.9 \\
Adult (ref.)   & 70  & 170 & 38.7 & 11.4 \\
Large adult    & 100 & 180 & 45.6 & 13.5 \\
\bottomrule
\end{tabular}
\end{table}

The reference adult clears the ICNIRP general-public reference level
of 10\,W/m$^2$ at 28\,GHz under worst-case directional exposure,
while smaller body sizes carry lower thresholds. The scaling
follows from the $m/A$ dependence in~\eqref{eq:Sinc-max-worst} and
matches the observation in the dosimetry literature~\cite{Hirata2007corr,
Dimbylow2002} that absorption cross-section scales with surface area
while mass scales with volume.

\subsection{Polarisation worst case}

Linearly polarised illumination at the most sensitive direction
shifts $\Tbar$ in~\eqref{eq:Sinc-max-worst} by at most a fraction
$|\tilde{B}|/(2\tilde{A}_\perp) \leq D_B$, with $D_B = 16\,\%$ on the
Thelonious anatomical phantom and $D_B \leq 27.9\,\%$ on the worst-
case convex shape (an infinite cylinder at broadside). For multipath
environments with $N \geq 20$ independent path orientations the
correction is suppressed below 2.5\,\%. The unpolarised
\eqref{eq:Sinc-max-worst} is therefore conservative for any
reverberation chamber and for the most common exposure scenario, a
vertically polarised base-station illumination of a standing person.


\section{Comparison with prior measurement campaigns}\label{sec:waterfall}

\Cref{eq:cauchy-exact} reproduces every plateau, dip, and cross-section
scaling reported in the reverberation-chamber and FDTD literature.
\Cref{fig:waterfall} collects the comparison.

\begin{figure*}[!t]
  \centering
  \includegraphics[width=\linewidth]{lit_waterfall.pdf}
  \caption{One closed-form identity matches every direction-averaged
  whole-body absorption ratio reported in the dosimetry literature
  above 6\,GHz, across 168 volunteers and 5 FDTD phantoms. (a)
  Flintoft 2014, 60 volunteers, 1--12\,GHz. The
  $\gamma_s = 0.865$-corrected points (filled triangles) collapse the
  5--11\,GHz plateau onto $\Tbar(f)$ within 2--4\,\%. (b) Zhang 2017,
  48 subjects, 1--18\,GHz. The prediction $\Tbar\cdot\Aab/A$ tracks the
  median in the geometric-optics range. (c) Bamba 2014, 4 phantoms
  FDTD, 1.45--5.8\,GHz, against $\Tbar(f)$. (d) Diao 2024
  (anatomical FDTD, 28\,GHz) and Kodera 2024 (parametric FDTD, 10 to
  100\,GHz) on $T_0(f)$. (e) Direct AEGIS / Sim4Life FDTD validation
  on the Thelonious phantom (Section~\ref{sec:fdtd}). (f) Common
  axis: every dataset converted to $\Tbar\cdot\Aab/A$ with one
  curve. The body-Mie / fat-resonance regime below 3\,GHz
  is shaded grey, where the layered tissue model in
  Section~\ref{sec:fp} recovers the dip.}
  \label{fig:waterfall}
\end{figure*}

\Cref{tab:waterfall} states the comparisons numerically. The Bamba
offset of 5--10\,\% in panel (c) is informative. Their $\eta$ is fit
from full-body FDTD on ellipsoidal phantoms in diffuse-field
exposure, and absorbs creeping-wave and finite-curvature
contributions that the planar-tissue $\Tbar$ omits. The convergence
of $\eta$ towards $\Tbar$ as frequency rises (and $ka$ increases)
is the convergence to the geometric-optics regime predicted by a Mie
analysis of body-scale spheres~\cite{BohrenHuffman1983}. At
5.8\,GHz, where Bamba's fit is at the upper edge of its calibration
range, the offset narrows to 5\,\%.

\begin{table*}[!t]
\centering
\caption{Quantitative comparison of the closed-form prediction to the
empirical literature. The shaded sub-3\,GHz regime, where the layered
tissue model in Section~\ref{sec:fp} replaces $\Tbar$, is excluded.}
\label{tab:waterfall}
\begin{tabular}{p{0.34\linewidth}p{0.34\linewidth}c}
\toprule
Empirical observation & Closed-form prediction & Match \\
\midrule
Flintoft~\cite{Flintoft2014}, $\langle Q^a\rangle/\gamma_s = 0.47$--$0.49$ at 7--11\,GHz, 60 volunteers
& $T_0(f) = 0.48$ at 9\,GHz, $\Tbar(f) = 0.50$
& 2--4\,\% \\
Bamba~\cite{Bamba2014}, $\eta(f) = 0.50$--$0.56$ at 1.5--5.8\,GHz, 4 FDTD phantoms
& $\Tbar(f) = 0.47$--$0.50$
& 5--10\,\% \\
Zhang~\cite{Zhang2017thesis}, $\xi = 0.45$--$0.65$ at 6--18\,GHz, 48 subjects
& $\Tbar(f)\cdot\Aab/A = 0.43$--$0.49$
& within scatter \\
Kodera~\cite{Kodera2024}, $T_{\mathrm{tr}}$ matches $A_\perp$ scaling at 10--100\,GHz, parametric FDTD
& $T_{\mathrm{tr}} \equiv T_0$
& $\leq 5\,\%$ \\
Diao~\cite{Diao2024}, $T = 0.52$ at 28\,GHz, anatomical FDTD
& $T_0 = 0.536$ at 28\,GHz
& 3\,\% \\
Flintoft, $-0.0061\,\mathrm{mm}^{-1}$ slope of $\langle Q^a\rangle$ vs $d_{\mathrm{SF}}$ at 3\,GHz
& Layered Fabry-P\'erot in fat
& mechanism, qualitative \\
\bottomrule
\end{tabular}
\end{table*}

Kodera et al.~\cite{Kodera2024} report the closest numerical
counterpart to the present analysis. Their Fig.~13 compiles WBASAR
data across roughly ten prior numerical and experimental studies
(TARO, HANAKO, Bahillo, Kuhn, Christ, Andersen, Hirata 2008,
Drossos, plus the four parametric Models I--IV) at 10\,W/m$^2$ from 1
to 100\,GHz, and shows the same asymptotic plateau that
\eqref{eq:cauchy-exact} predicts. They fit a study-specific
$T_{\mathrm{tr}}$ per phantom and frequency.
\Cref{eq:cauchy-exact} supplies the closed-form $T_{\mathrm{tr}} \to
\Tbar(f)$ that all phantoms converge to in the geometric-optics
regime, with the residual phantom-to-phantom spread set by the
body-shape factor $\Aab/A$.


\section{FDTD validation on the Thelonious phantom}\label{sec:fdtd}

The literature comparison shows consistency with three
independent volunteer studies and two FDTD studies. A single-
laboratory FDTD ground-truth comparison on the Thelonious anatomical
phantom completes the validation. The setup uses full Sim4Life FDTD
on the Thelonious mesh (23\,826 triangles, 0.787\,m$^2$ surface area,
17.4\,kg), with twelve incidence directions, two polarisations, and
nine frequencies from 450\,MHz to 5.8\,GHz. The prediction is
\eqref{eq:cauchy-exact} with $\Aab/A = 0.865$ from the AEGIS
ambient-occlusion solver and $\Tbar(f)$ from \cref{tab:Tbar}, with
no fitted parameters.

\Cref{fig:val-fdtd} shows the kernel-by-kernel ratio of the AEGIS
prediction to the FDTD result, direction-averaged over $12 \times 2 =
24$ scenarios per frequency. The closed-form Cauchy prediction (black
stars in the right panel) approaches unity at the upper end of the
band, with the ratio reaching 1.012 at 5.8\,GHz. Below 5\,GHz the
prediction underestimates because body-scale Mie and resonance effects
do not enter a surface-only law. The breakdown follows the Mie
scaling on a sphere of comparable size parameter. The layered-tissue
extension in Section~\ref{sec:fp} recovers the sub-5\,GHz behaviour
qualitatively but not quantitatively at the body-averaged level.

\begin{figure*}[!t]
  \centering
  \includegraphics[width=\linewidth]{fig_kernels_vs_fdtd.png}
  \caption{Direction-averaged AEGIS / FDTD comparison on the
  Thelonious phantom from 0.45 to 5.8\,GHz. Twelve incidence
  directions and two polarisations per frequency. The closed-form
  Cauchy prediction (black stars, right panel) reaches 1.012 at
  5.8\,GHz with $\Aab/A = 0.865$ and $\Tbar(f)$ both computed in
  AEGIS with no fitted parameters.}
  \label{fig:val-fdtd}
\end{figure*}

The non-convexity factor $\Aab/A = 0.865$ used in this comparison is
computed by AEGIS from the same Thelonious mesh with no calibration.
A separate calculation using the cosine-weighted hemispheric integral
on the same mesh gives $\Aab/A = 0.87$. The two independent
computations agree to 0.6\,\%, an internal consistency check on the
geometric pipeline. Flintoft's geometric estimate from Tomita's
surface-area data lies in 0.75--0.85~\cite{Flintoft2014}. The present
value sits at the upper end of his band, consistent with Thelonious
being a six-year-old child phantom (smaller torso-arm gap relative to
adults).


\section{Reverberation-chamber dosimetry}\label{sec:reverb}

A reverberation chamber produces a statistically isotropic and
uniformly polarised electromagnetic field. The direction average
of~\eqref{eq:cauchy-exact} applies, and the absorbed power on the
body in the chamber is
\begin{equation}\label{eq:reverb}
  P_{\mathrm{abs}} = S_{\mathrm{total}}\,\Tbar(f)\,\Aab/4,
\end{equation}
where $S_{\mathrm{total}}$ is the chamber's measured scalar power
density. Three precomputed scalars determine the result: $\Tbar$ from
one Fresnel quadrature, $\Aab$ from one ambient-occlusion pass on the
body mesh, and the body mass $m$ for the SAR conversion. No direction
sampling, no FDTD per configuration, no calibration step between
phantom and chamber.

This is the regime where the closed form matters most for the
dosimetry community. Reverberation-chamber experiments on small
animals (mice, rats), cell cultures, and human-body phantoms have for
two decades required FDTD on the specific phantom-chamber combination
to estimate dose, calibrated against measured $\sigma_a$ via Bamba's
empirical $\eta$~\cite{Bamba2014} or Flintoft's
$\langle Q^a\rangle$~\cite{Flintoft2014}. \Cref{eq:reverb} replaces
the FDTD step with one closed-form evaluation. For small animals with
$\eta \approx 1$ (convex morphology), the formula reduces to one
Fresnel quadrature and one mass measurement.


\section{Discussion}\label{sec:disc}

The closed-form identity~\eqref{eq:cauchy-exact} carries three
constraints on its applicability. First, the body must be
electromagnetically opaque, which holds above $\sim 100\,$MHz.
Second, the illumination must be statistically isotropic for the
unpolarised direction average to apply. The polarisation correction
in Section~\ref{sec:compliance} bounds the worst case. Third,
body-scale Mie and diffraction effects below 5\,GHz reduce the
identity to a qualitative tool in that band. The layered-tissue
extension in Section~\ref{sec:fp} accounts for the dominant
fat-layer Fabry-P\'erot resonance but not the broader Mie behaviour.

A 20\,\% uncertainty floor enters from the tissue dielectric tables
themselves~\cite{AlekseevZiskin2007}. The identity sits within this
floor against Sim4Life FDTD on the Thelonious phantom (1.2\,\% at
5.8\,GHz), against every plateau and dip in 168 volunteers across the
Bamba 2014, Flintoft 2014, and Zhang 2017 campaigns, and against the
parametric and anatomical FDTD campaigns of Diao 2024 and Kodera 2024
above 10\,GHz.

Identifying the exposure fraction $\eta$ with the ambient-
occlusion factor of computer graphics has a practical implication.
Four decades of GPU-optimised algorithms apply directly to whole-
body dosimetry. For a typical anatomical mesh of $10^4$--$10^5$
triangles, $\eta(\rr)$ evaluates in tens of milliseconds on commodity
hardware, against minutes to hours for the FDTD step it replaces.

\section{Conclusion}\label{sec:concl}

The direction-averaged whole-body absorbed power on any opaque body in
a statistically isotropic field equals $\Sinc\,\Tbar(f)\,\Aab/4$. The
flux-weighted transmission $\Tbar(f)$ is one quadrature against the
IT'IS dielectric tables. The absorption area $\Aab$ is one ambient-
occlusion pass over the body mesh. For reverberation-chamber
dosimetry of small animals and human-body
phantoms,~\eqref{eq:reverb} replaces the FDTD calibration step the
community has used for two decades. For ICNIRP whole-body SAR
compliance,~\eqref{eq:Sinc-max-worst} reduces the threshold to a
closed-form function of three precomputed scalars: $\Tbar$, $\Aab$,
and the body mass $m$.


\bibliographystyle{IEEEtran}
\bibliography{refs}

\end{document}
